\end{itemize}
We call $T$ a {\it partial tilting object} if the conditions (1) and (2) are satisfied.
\end{dfn}

For an object $T\in \scrA$, we denote by
\begin{gather*}
\Add T
\end{gather*}
the smallest additive subcategory containing all direct summands of small coproducts of $T$. We~write $\add T\subset \Add T$ for the subcategory consisting of direct summands of finite coproducts of~$T$. If~$T$ is a partial tilting object, then for arbitrary sets $\cI$ and $\cJ$, we have isomorphisms
\begin{align*}
\Ext^n_{\scrA}\Big(\bigoplus_{i\in \cI}T,\bigoplus_{j\in \cJ}T\Big) &\cong\Hom_{\D(\scrA)}\Big(\bigoplus_{i\in \cI}T,\bigoplus_{j\in \cJ}(T[n])\Big)
\\
&\cong \prod_{i\in \cI}\Hom_{\D(\scrA)}\Big(T,\bigoplus_{j\in \cJ}(T[n])\Big)
\\
&\cong \prod_{i\in\cI}\Big(\bigoplus_{j\in\cJ}\Hom_{\D(\scrA)}(T,T[n])\Big).
\end{align*}
This implies that for any $X,Y\in \Add T$, we have $\Ext_{\scrA}^n(X,Y)=0$ for any $n>0$. Thus any short exact sequence $0\to X\to Z\to Y\to0$ in $\scrA$ with $X,Y\in \Add T$ splits, and in particular $\Add T$ is an exact subcategory of $\scrA$.

In the remaining of this subsection, $T\in\scrA$ is a partial tilting object, and we assume that every right module over the endomorphism ring $\Lambda\defeq\End_{\scrA}(T)$ is of finite projective dimension. Consider the functor
 \begin{gather*}
{\F}\defeq\Hom_{\scrA}(T,-)\colon\ \scrA\to\Mod\Lambda
 \end{gather*}
and assume that
\begin{itemize}\itemsep=0pt
 \item[$(i)$] $\F$ has a left adjoint functor $\G\colon\Mod\Lambda\to \scrA$ such that the adjunction morphisms $\sigma({\Lambda})\colon \Lambda\to \F\G(\Lambda)$ and $\tau(T)\colon \G\F(T)\to T$ are isomorphisms.

 \item[$(ii)$] The right derived functor $\bR \F\colon \D^+(\scrA)\to \D^+(\Mod\Lambda)$ restricts to the functor
\begin{gather*}
 \bR \F\colon\ \Db(\scrA)\to \Db(\Mod\Lambda)
 \end{gather*}
 between bounded derived categories.
 \end{itemize}

 Note that both of the functors $\F$ and $\G$ commute with small coproducts since $T$ is compact and $\G$ admits a right adjoint functor. In particular, $\F$ and $\G$ restrict to the following functors
 \begin{gather*}
 \F\colon\ \Add T\to\Proj\Lambda,
 \\
\G\colon\ \Proj\Lambda\to \Add T,
 \end{gather*}
where $\Proj \Lambda$ is the category of projective right $\Lambda$-modules.
\begin{lem}\label{add proj}
The functor $\F\colon\Add T\to\Proj\Lambda$ is an equivalence, and $\G\cong \F^{-1}$.
\begin{proof}
The adjunction morphisms $\sigma({P})\colon P\to \F\G(P)$ and $\tau(S)\colon \G\F(S)\to S$ are bijective for any $P\in\Mod\Lambda$ and $S\in \Add T$, since $\sigma({\Lambda})$ and $\tau({T})$ are bijective and $\Proj \Lambda=\Add \Lambda$. Hence, the functors $\F$ and $\G$ are equivalences, and $\G\cong \F^{-1}$.
\end{proof}
\end{lem}

By our assumptions, the functor $\F$ defines the right derived functor
\begin{gather*}
\bR\F\colon\ \Db(\scrA)\to \Db(\Mod\Lambda)
\end{gather*}
between bounded derived categories, and the functor $\G$ also defines the left derived functor
\begin{gather*}
\bL\G\colon\ \Db(\Mod\Lambda)\to\Db(\scrA)
\end{gather*}
between bounded derived categories.
The following is well known to experts.


\begin{thm}\label{general tilting}
The functor $\bL \G\colon \Db(\Mod\Lambda)\to\Db(\scrA)$
is fully faithful. Furthermore, if $T$ is a tilting object, then $\bL \G$ is an equivalence and $\bR \F\cong (\bL \G)^{-1}$.
\begin{proof}
 We have the following commutative diagram
\[
\begin{tikzcd}
\Kb(\Proj\Lambda)\arrow[rr,"\G"]\arrow[d,]&&\Kb(\Add T)\arrow[d,]\\
\Db(\Mod\Lambda)\arrow[rr, "\bL \G\,\,\,\,\,\,"]&&\Db(\scrA),
\end{tikzcd}
\]
where the vertical arrows are the natural quotient functors. The left vertical functor is an equi\-va\-lence, since every right $\Lambda$-module has a finite projective resolution. Moreover, by Lemma~\ref{add proj}, the top horizontal functor is also an equivalence. Hence, for the former assertion, it is enough to show that the natural functor
\begin{gather*}
\Kb(\Add T)\to \Db(\scrA)
\end{gather*}
is fully faithful, and this follows from an identical argument as in~\cite[Chapter~III, Lemma~2.1]{happel}.

For the latter assertion, assume that $T$ is a tilting object. It is enough to prove that $\bR \F$ is also fully faithful. For any $A^{\bullet}\in \D(\scrA)$, consider the following triangle
\begin{gather*}
\bL \G \bR \F(A^{\bullet})\xrightarrow{\tau}A^{\bullet}\to C(\tau)\to \bL \G \bR \F(A^{\bullet})[1],
\end{gather*}
where $\tau$ is the adjunction morphism. Applying the functor $\bR \F$ to the above triangle and using the natural isomorphism $\bR \F\circ \,\bL \G\cong \id$, we see that $\bR \F(C(\tau))\cong 0$ in $\Db(\Mod\Lambda)$. Since $T$ is a~generator of $\D(\scrA)$, this implies that $C(\tau)\cong 0$ in $\D(\scrA)$. Hence the adjunction morphism $\tau$ is an~isomorphism. This completes the proof.
\end{proof}
\end{thm}

\begin{rem}
Our assumption that every $\Lambda$-module has a finite projective resolution might be weakened. However, the result under our assumption is enough for our purpose.
\end{rem}

\section{Equivariant algebras and equivariant tilting modules}

In this section, we introduce equivariant algebras and equivariant modules, and define seve\-ral functors between equivariant modules that are generalizations of functors in~\cite[Section~2]{bfk}. We~also show that equivariant tilting modules induce derived equivalences of equivariant modules, which are simultaneous generalizations of tilting equivalences induced by tilting bundles on schemes and tilting modules over noncommutative algebras.

\subsection{Equivariant modules over equivariant algebras}

In this subsection, we give the definitions of equivariant algebras and equivariant modules, and~dis\-cuss basic properties.
Notation is the same as in Section~\ref{equiv section}. Let $\uprho\colon \cO_X\to \cA$ be an $X$-algebra.

\begin{dfn}
A {\it $G$-equivariant structure} on the $X$-algebra $\uprho\colon \cO_X\to\cA$ is an isomorphism
\begin{gather*}
\uptheta^{\cA}\colon\ \pi^*\cA\simto \sigma^*\cA
\end{gather*}
 of sheaves of algebras on $G\times X$
such that it is $\cO_{G\times X}$-linear and it gives a $G$-equivariant structure on the $\cO_X$-module $\cA$. The pair $\big(\cA,\uptheta^{\cA}\big)$ is called a {\it $G$-equivariant algebra over $X$}, or {\it $G$-equivariant $X$-algebra}. We say that $\big(\cA,\uptheta^{\cA}\big)$ is {\it coherent} if the underlying $X$-algebra $\cA$ is coherent $\cO_X$-module. We write simply $\cA$ for $\big(\cA,\uptheta^{\cA}\big)$ if no confusion seems likely to occur.
\end{dfn}

\begin{rem}\label{linearity}
Note that for an isomorphism $\upphi \colon \pi^*\cA\simto \sigma^*\cA$ of sheaves of algebras which defines a $G$-equivariant structure on the quasi-coherent sheaf $\cA$, $\upphi$ is $\cO_{G\times X}$-linear if and only if the morphism $\uprho\colon \cO_X\to \cA$ is a $G$-equivariant morphism of $G$-equivariant quasi-coherent sheaves $(\cO_X,\uptheta^{\rm can})$ and $(\cA,\upphi)$, where $\uptheta^{\rm can}$ is the canonical equivariant structure on the structure sheaf~$\cO_X$ induced by the $G$-action on $X$.
\end{rem}

Let $X=\Spec R$ be an affine scheme with an action from an affine algebraic group $G$, and set $R_G\defeq k[G]\otimes_kR$, where $k[G]$ denotes the coordinate ring of $G$. We denote by
\begin{gather*}
R_G^{\pi}\qquad \big(\text{resp.}~R_G^{\sigma}\big)
\end{gather*}
the $R$-algebra $R\to R_G$ induced by the morphism $\pi\colon G\times X\to X$ (resp.\ $\sigma\colon G\times X\to X$).
If~$\big(\cA,\uptheta^{\cA}\big)$ is a $G$-equivariant $X$-algebra, taking global sections induces an $R$-algebra $A\defeq\cA(X)$ and an isomorphism
\begin{gather*}
\uptheta^{A}\defeq\uptheta^{\cA}(X)\colon\ A\otimes_R R_G^{\pi}\simto A\otimes_RR_G^{\sigma}
\end{gather*}
of $R_G$-algebras such that $\uptheta^{A}$ gives a $G$-equivariant structure on the $R$-module $A$. This yields the following definition.


 \begin{dfn}
 We call a pair
 $(\Lambda,\uptheta^{\Lambda})$
 a {\it $G$-equivariant $R$-algebra}, if $\Lambda$ is an $R$-algebra and $\uptheta^{\Lambda}\colon \Lambda\otimes_RR_G^{\pi}\simto\Lambda\otimes_RR_G^{\sigma}$ is an isomorphism of $R_G$-algebras such that it defines a $G$-equivariant structure on the $R$-module $\Lambda$.
 \end{dfn}

 \begin{rem}
For any commutative ring $S$, via the natural equivalence $\Mod S\simto \Qcoh \Spec S$, to give a $G$-equivariant algebra over $\Spec S$ is equivalent to give a $G$-equivariant $S$-algebra.
 \end{rem}

\begin{exa}\label{end ex}
Notation is the same as above.
\begin{enumerate}\itemsep=0pt
\item Let $\cF\in \coh_GX$ be a $G$-equivariant coherent sheaf on $X$. Then the algebra
%\label{end ex 1}
\begin{gather*}
\cA\defeq\cEnd_X(\cF)
\end{gather*}
over $X$ has a natural $G$-equivariant structure. Indeed, the $G$-equivariant structure
\begin{gather*}
\uptheta\colon\ \pi^*\cEnd_X(\cF)\simto \sigma^*\cEnd_X(\cF)
\end{gather*}
on the $G$-equivariant coherent sheaf $\cEnd_X(\cF)\in\coh_GX$
is a morphism of algebras, and so the pair $(\cA, \uptheta$) is a $G$-equivariant algebra over $X$.


\item Let $\big(\cA,\uptheta^{\cA}\big)$ be a $G$-equivariant $X$-algebra. If $Y$ is another $G$-scheme and $f\colon X\to Y$ is a~$G$-equ\-i\-variant morphism that is quasi-separated and quasi-compact,
then the direct image
$\big(f_*\cA,(\id_G\times f)_*\uptheta^{\cA}\big)$
is a $G$-equivariant $Y$-algebra. If $\cA$ is coherent and $f$ is proper, the algebra $f_*\cA$ is also coherent. Similarly, if $g\colon Y\to X$ is a $G$-equivariant morphism, then the pull-back $\big(g^*\cA,(\id_G\times g)^*\uptheta^{\cA}\big)$ is a $G$-equivariant $Y$-algebra, and it is coherent if~$\cA$ is coherent.
%\label{end ex 2}
 \end{enumerate}
\end{exa}



We define equivariant modules over equivariant algebras.

\begin{dfn}\label{equiv sheaf}
Let $\big(\cA,\uptheta^{\cA}\big)$ be a $G$-equivariant $X$-algebra.
\begin{itemize}\itemsep=0pt
\item[1.] A {\it $G$-equivariant structure} on a quasi-coherent right $\cA$-module $\cM$ is an isomorphism
\begin{gather*}
\uptheta^{\cM}\colon\ \pi^*\cM\simto \sigma^*\cM
\end{gather*}
of $\pi^*\cA$-modules such that $\uptheta^{\cM}$ satisfies the condition \eqref{compati}, where the $\pi^*\cA$-module structure on the $\sigma^*\cA$-module $\sigma^*\cM$ is given by $\uptheta^{\cA}$.
We call the pair $\big(\cM,\uptheta^{\cM}\big)$ a {\it $G$-equivariant quasi-coherent right $\cA$-module}, or simply {\it $G$-equivariant $\cA$-module}. If the underlying sheaf~$\cM$ of~modu\-les is coherent, we call $\big(\cM,\uptheta^{\cM}\big)$ a {\it $G$-equivariant coherent $\cA$-module}. We~sometimes write just $\cM$ for $\big(\cM,\uptheta^{\cM}\big)$.


\item[2.] Let $\big(\cM,\uptheta^{\cM}\big)$ and $\big(\cN,\uptheta^{\cN}\big)$ be $G$-equivariant $\cA$-modules. A morphism $\varphi\colon \cM\to \cN$ of $\cA$-modules is {\it $G$-equivariant} if the following diagram
\[
\begin{tikzcd}
\pi^*\cM\arrow[rr,"\pi^*\varphi"]\arrow[d,"\uptheta^{\cM}"']&&\pi^*\cN\arrow[d," \uptheta^{\cN}"]\\
\sigma^*\cM\arrow[rr, "\sigma^*\varphi"]&&\sigma^*\cN
\end{tikzcd}
\]
is commutative.

\item[3.] We denote by
\begin{gather*}\Qcoh_G\cA\end{gather*} the category of $G$-equivariant $\cA$-modules whose morphisms are $G$-equivariant, and write $\coh_G\cA$ for the full subcategory of $G$-equivariant coherent $\cA$-modules.
\end{itemize}
\end{dfn}

\begin{rem}
(1) If $X$ is a $G$-scheme, then $\cO_X$ is a natural $G$-equivariant algebra over $X$. For any $(\cF,\uptheta)\in \Qcoh_GX$, the $G$-equivariant structure $\uptheta$ is automatically $\pi^*\cO_X$-linear. Therefore, we have a natural identification
\begin{gather*}
\Qcoh_G\cO_X=\Qcoh_GX.
\end{gather*}
Hence $G$-equivariant modules are generalizations of $G$-equivariant quasi-coherent sheaves.
\end{rem}

\begin{dfn}%\label{equiv module}
Let $X=\Spec R$ be an affine scheme with an action from an affine algebraic group $G$, and $(\Lambda,\uptheta^{\Lambda})$ a $G$-equivariant $R$-algebra. A {\it $G$-equivariant structure} on a right $\Lambda$-module~$M$ is an isomorphism
\begin{gather*}
\uptheta^M\colon\ M\otimes_RR_G^{\pi}\simto M\otimes_RR_G^{\sigma}
\end{gather*}
of $\big(\Lambda\otimes_RR_G^{\pi}\big)$-modules, where $\big(\Lambda\otimes_RR_G^{\sigma}\big)$-module $M\otimes_RR_G^{\sigma}$ is considered as a $\big(\Lambda\otimes_RR_G^{\pi}\big)$-mo\-dule via the ring isomorphism $\uptheta^{\Lambda}\colon \Lambda\otimes_RR_G^{\pi}\simto\Lambda\otimes_RR_G^{\sigma}$. We call such a pair $\big(M,\uptheta^M\big)$ a~{\it $G$-equivariant $\Lambda$-module}. We denote by
\begin{gather*}
\Mod_G\Lambda
\end{gather*}
the category of $G$-equivariant $\Lambda$-modules whose morphisms are defined similarly to Definition~\ref{equiv sheaf}, and we denote by
\begin{gather*}
\fmod_G\Lambda\subset \Mod_G\Lambda
\end{gather*}
the full subcategory consisting of equivariant modules that are finitely generated over $\Lambda$.
\end{dfn}

\begin{rem}
Let $X=\Spec R$ be an affine scheme with an action from an affine algebraic group $G$, and $\big(\cA,\uptheta^{\cA}\big)$ a $G$-equivariant $X$-algebra. Then it induces a $G$-equivariant $R$-algebra $A\defeq\cA(X)$, and we have a natural equivalence
\begin{gather*}
\Qcoh_G\cA\simto\Mod_GA
\end{gather*}
of abelian categories, which restricts to an equivalence $\coh_G\cA\simto\fmod_GA$.
\end{rem}
\vspace{1mm}


Let $\big(\cM,\uptheta^{\cM}\big)$ be a $G$-equivariant $\cA$-module. Recall from Remark~\ref{closed point isom} that for each closed point $g\in G$, we have the induced isomorphism
\begin{gather*}
\uptheta_g^{\cM}\colon\ \cM\simto \sigma_g^*\cM.
\end{gather*}
If $\big(\cN,\uptheta^{\cN}\big)$ is another $G$-equivariant $\cA$-module, we have the $G$-action on $\Hom_{\Qcoh\cA}(\cM,\cN)$ defined by
\begin{gather*}
g\cdot\varphi\defeq\big(\uptheta^{\cN}_g\big)^{-1}\circ\sigma_g^*\varphi\circ\uptheta^{\cM}_g.
\end{gather*}
The following is a generalization of the equality in \eqref{G hom}.
\begin{prop}\label{G hom 2}
Assume that $X$ is of finite type over $k$. We have
\begin{gather*}
\Hom_{\Qcoh_G\cA}\big(\big(\cM,\uptheta^{\cM}\big),\big(\cN,\uptheta^{\cN}\big)\big) =\Hom_{\Qcoh\cA}(\cM,\cN)^G.
\end{gather*}
Moreover, if $G$ is reductive, we have
\begin{gather*}
\Ext_{\Qcoh_G\cA}^i\big(\big(\cM,\uptheta^{\cM}\big),\big(\cN,\uptheta^{\cN}\big)\big) =\Ext_{\Qcoh\cA}^i(\cM,\cN)^G.
\end{gather*}
for any $i\geq0$.
\begin{proof}
The first equality follows from an identical argument as in Remark~\ref{cl pt}. Since $G$ is reductive, the functor
of taking $G$-invariant parts is exact, and thus the second equality follows from the first one.
\end{proof}
\end{prop}

\begin{exa} %\label{hom tens}
Notation is the same as in Example~\ref{end ex}(1).

\begin{itemize}\itemsep=0pt
\item[$1.$] For any $\cG\in \Qcoh_GX$, the quasi-coherent $\cA$-module \begin{gather*}\cHom_X(\cF,\cG)\end{gather*} has a natural $G$-equivariant structure induced by the $G$-equivariant structure on
 \begin{gather*}
 \cHom_X(\cF,\cG)\in \Qcoh_GX.
 \end{gather*}


\item[$2.$] For $\cM\in \Mod_G\cA$ and $\cF\in \Qcoh_GX$, the tensor product \begin{gather*}\cF\otimes_X \cM\end{gather*} of $\cO_X$-modules is an $\cA$-module with a natural $G$-equivariant structure induced by the $G$-equivariant structure on $\cF\otimes_X \cM\in \Qcoh_GX$.
\end{itemize}
\end{exa}

If $\big(\cA,\uptheta^{\cA}\big)$ is a $G$-equivariant $X$-algebra, the ring homomorphism
$\uprho\colon \cO_X\to \cA$ is $G$-equivariant by Remark~\ref{linearity}. We define the sheaf of rings
 \begin{gather*}
 [\cA/G]
 \end{gather*}
 on the quotient stack $[X/G]$ to be the image of $\cA\in \Qcoh_GX$ by the equivalence
\begin{gather*}
[-/G]\colon\ \Qcoh_GX\simto \Qcoh[X/G]
\end{gather*}
 in \eqref{equiv stack}. Since $[\cO_X/G]$ is canonically isomorphic to the structure sheaf $\cO_{[X/G]}$, we have the following morphism
\begin{gather*}
[\uprho/G]\colon\ \cO_{[X/G]}\to [\cA/G]
\end{gather*}
of sheaf of rings and this makes the sheaf $[\cA/G]$ an algebra over $[X/G]$.
The following is a~generalization of \eqref{equiv stack}, and it follows from Proposition~\ref{main app}.
\begin{prop}\label{quotient correspond}
We have an equivalence
\begin{gather*}
\Qcoh_G\cA\cong \Qcoh [\cA/G]
\end{gather*}
of abelian categories.
\end{prop}

\begin{rem}
If an affine algebraic group $G$ is abelian, by a similar argument as in~\cite[Section~2.1]{bfk2} $G$-equivariant algebras correspond to $\widehat{G}$-graded algebras, where $\widehat{G}$ is the character group of $G$. Since we do not need this correspondence in the present paper, we do not give a~formulation of the correspondence and its proof.
\end{rem}

\subsection{Functors of equivariant modules}
In this subsection, we define fundamental functors between equivariant modules. Let $X$ be a~qua\-si-compact and quasi-separated scheme, and $G$ an algebraic group acting on $X$ by $\sigma\colon G\times X\to X$. Denote by $\pi\colon G\times X\to X$ the natural projection.

\subsubsection{Restrictions and extensions}

Let $\cA$ and $\cB$ be $G$-equivariant $X$-algebras, and let $\varphi\colon \cA\to \cB$ be a $G$-equivariant morphism of~$X$-algebras. The functors in \eqref{restriction} and \eqref{extension} define the {\it restriction} by $\varphi$
\begin{gather*}
(-)_{\varphi}\colon\ \Qcoh_G\cB\to \Qcoh_G\cA
\end{gather*}
and the {\it extension} by $\varphi$
\begin{gather*}
(-)\otimes_{\cA}\cB\colon\ \Qcoh_G\cA\to \Qcoh_G\cB,
\end{gather*}
and we have an adjunction
\begin{equation}\label{rest ext adj}
(-)\otimes_{\cA}\cB\dashv (-)_{\varphi}.
\end{equation}

\subsubsection{Direct image functors and pull-back functors}%\label{3.2.1}
Let $Y$ be another quasi-compact and quasi-separated $G$-scheme, and $f\colon X\to Y$ a $G$-equivariant morphism.
Let $\cA$ be a $G$-equivariant $X$-algebra, and $\cB$ a $G$-equivariant $Y$-algebra. Recall from Example~\ref{end ex}(2) that the push-forward $f_*\cA$ is a $G$-equivariant $Y$-algebra and that the pull-back $f^*\cB$ is a $G$-equivariant $X$-algebra. The push-forwards and the pull-backs of $G$-equivariant modules define the following additive functors
\begin{gather*}
f_{\star}\colon\ \Qcoh_G\cA\to \Qcoh_Gf_*\cA,
\\[.5ex]
f^*\colon\ \Qcoh_G\cB \to \Qcoh_Gf^*\cB.
\end{gather*}
Using natural morphisms of equivariant algebras
$\varphi\colon f^{-1}f_*\cA\to \cA$ and $\psi\colon \cB\to f_*f^*\cB$,
we define additive functors
\begin{gather*}
f^{\star}\colon\ \Qcoh_Gf_*\cA\to\Qcoh_G\cA,
\\[.5ex]
f_{*}\colon\ \Qcoh_Gf^*\cB\to \Qcoh_G\cB
\end{gather*}
by $f^{\star}(-)\defeq f^{-1}(-)\otimes_{f^{-1}f_*\cA}\cA$ and $f_{*}\defeq(-)_{\psi}\circ f_{\star}$.
By standard arguments, we see that the above functors induce adjunctions;
\begin{gather*}
f^*\dashv f_{*}\qquad\text{and}\qquad f^{\star}\dashv f_{\star}.
\end{gather*}




 Let $G'$ be another algebraic group acting on another scheme $Z$, and let $\alpha\colon G\to G'$ be a morphism of algebraic groups. Let $h\colon X\to Z$ be a morphism of schemes, and $\cC$ an $G'$-equivariant $Z$-algebra. We say that $h$ is {\it $\alpha$-equivariant} if we have $h(g x)=\alpha(g)h(x)$ for all $(g,x)\in G\times X$. If $h$ is $\alpha$-equivariant, then $h$ induces the pull-back functor
\begin{gather*}
h^*_{\alpha}\colon\ \Qcoh_{G'}Z\to \Qcoh_GX
\end{gather*}
defined by $h^*(\cF,\uptheta)\defeq(h^*\cF, (\alpha\times h)^*\uptheta)$, and we have
 the associated $G$-equivariant $X$-alge\-bra~$h^*\cC$. This functor extends to
the functor
\begin{equation*}%\label{relative pull-back}
h^*_{\alpha}\colon\ \Qcoh_{G'}\cC\to \Qcoh_Gh^*\cC.
\end{equation*}

\subsubsection{Tensor products and sheaf Homs}

Let $\uprho\colon\cO_{X}\to\cA$ and $\uprho'\colon\cO_{X}\to\cA'$ be $G$-equivariant $X$-algebras. A {\it $G$-equivariant quasi-coherent $($resp.\ coherent$)$ $(\cA',\cA)$-bimodule} is a pair $\big(\cM,\uptheta^{\cM}\big)$ of a $(\cA',\cA)$-bimodule $\cM$ such that the left action of $\cO_{X}$ via $\uprho'$ coincides with the right action of $\cO_{X}$ via $\uprho$ and that $\cM$ is a quasi-coherent (resp.\ coherent) sheaf on $X$ and an isomorphism $\uptheta\colon \pi^*\cM\simto\sigma^*\cM$ of $(\pi^*\cA',\pi^*\cA)$-bimodules such that $\uptheta^{\cM}$ satisfies the condition \eqref{compati}.
 We denote by
$\Qcoh_G(\cA',\cA)$ the category of $G$-equivariant quasi-coherent $(\cA',\cA)$-bimodules, and $\coh_G(\cA',\cA)$ the full subcategory of $G$-equivariant coherent $(\cA',\cA)$-bimodules.


Let $\cF\in \Qcoh_G(\cA',\cA)$, $\cM\in\Qcoh_G\cA'$ and $\cN\in \Qcoh_G\cA$. Then we have the tensor product
\begin{gather*}
\cM\otimes_{\cA'}\cF\in \Qcoh_G\cA.
\end{gather*}
If $\cF\in \coh_G(\cA',\cA)$, we also have the sheaf Hom
\begin{gather*}
\cHom_{\cA}(\cF,\cN)\in \Qcoh_G\cA',
\end{gather*}
and there is a functorial isomorphism
\begin{equation*}%\label{adjoint tens}
\Hom_{\Qcoh_G\cA}\big(\cM\otimes_{\cA'}\cF,\cN\big)\cong \Hom_{\Qcoh_G\cA'}\big(\cM,\cHom_{\cA}(\cF,\cN)\big).
\end{equation*}
In other words, if $\cF\in \coh_G(\cA',\cA)$, the functor
\begin{gather*}
(-)\otimes_{\cA'}\cF\colon\ \Qcoh_G\cA'\to\Qcoh_G\cA
\end{gather*}
is left adjoint to the functor
\begin{gather*}
\cHom_{\cA}(\cF,-)\colon\ \Qcoh_G\cA\to\Qcoh_G \cA'.
\end{gather*}

The following is a version of projection formula for equivariant $\cA$-modules.

\begin{lem}[Projection formula]\label{proj form}
Let $Y$ be a quasi-compact and quasi-separated scheme with a~$G$-action, and $g\colon Y\to X$ a $G$-equivariant morphism. If $g$ is flat and affine, for any $\cE\in \Qcoh_G Y$ and $\cF\in \Qcoh_G\cA$, we have an isomorphism of $G$-equivariant $\cA$-modules
\begin{gather*}
g_*(\cE)\otimes_{\cO_X} \cF\simto g_*\!\big(\cE\otimes_{\cO_Y}g^*\cF \big).
\end{gather*}
\begin{proof}
Since $\bR g_*\cong g_*$ and $\bL \,g^*\cong g^*$, by the projection formula~\cite[Proposition~3.9.4]{lip}, we~have a quasi-isomorphism
\begin{gather*}
g_*\cE\otimes^{\bL}_{\cO_X}\cF\cong g_*\big(\cE\otimes^{\bL}_{\cO_Y}g^*\cF\big)
\end{gather*}
of complexes of quasi-coherent $\cO_X$-modules. Hence we have isomorphisms
\begin{align*}
g_*(\cE)\otimes_{\cO_X} \cF&\cong H^0\big(g_*(\cE)\otimes^{\bL}_{\cO_X}\cF\big)\\
&\cong H^0\big(g_*\big(\cE\otimes^{\bL}_{\cO_Y}g^*\cF\big)\big)\\
&\cong g_*\big(H^0\big(\cE\otimes^{\bL}_{\cO_Y}g^*\cF\big)\big)\\
&\cong g_*\big(\cE\otimes_{\cO_Y}g^*\cF\big),
\end{align*}
where the third isomorphism follows since $g_*$ is an exact functor. This isomorphism
\begin{gather*}
\varphi\colon\ g_*(\cE)\otimes_{\cO_X} \cF\simto g_*\!\big(\cE\otimes_{\cO_Y}g^*\cF \big)
\end{gather*}
of $\cO_X$-modules is nothing but the composition
\begin{gather*}
 g_*(\cE)\otimes_{\cO_X} \cF\longrightarrow g_*(\cE)\otimes_{\cO_X} g_*g^*\cF\longrightarrow g_*\!\big(\cE\otimes_{\cO_Y}g^*\cF \big),
\end{gather*}
where the first morphism is the morphism induced by the adjunction $\cF\to g_*g^*\cF$ and the second morphism is the canonical morphism defined by $x\otimes y\mapsto x\otimes y$. Since these morphisms are $\cA$-linear and $G$-equivariant, so is the composition $\varphi$. This finishes the proof.
\end{proof}
\end{lem}

\subsubsection{Taking invariant sections.}
Let $H$ be a closed normal subgroup of $G$. Assume that the restriction $\sigma_H\defeq\sigma|_{H\times X}\colon H\times X\to X$ is the trivial action, so that $\sigma_H=\pi_H$, where $\pi_H\defeq\pi|_{H\times X}\colon H\times X\to X$ is the natural projection. Then we have the induced $G/H$-action on $X$ denoted by
\begin{gather*}
\overline{\sigma}\colon\ G/H\times X\to X.
\end{gather*}
We write $\overline{\pi}\colon G/H\times X\to X$ for the natural projection.

For $\big(\cF,\uptheta^{\cF}\big)\in\Qcoh_GX$, we define the subsheaf $\cF^H\subseteq\cF$ to be the kernel of the composition
\[
\begin{tikzcd}
\cF\arrow[r,"\mu"]&(\pi_H)_*(\pi_H)^* \cF\arrow[rr,"\id -(\pi_H)_*\uptheta_H^{\cF}"]&&(\pi_H)_*(\pi_H)^* \cF,
\end{tikzcd}
\]
 where $\mu$ is the adjunction morphism and $\uptheta^{\cF}_H\defeq\uptheta^{\cF}|_{H\times X}\colon (\pi_H)^*\cF\simto (\sigma_H)^*\cF=(\pi_H)^*\cF$.
Then the pair $\big(\cF^H,\uptheta^{\cF}|_{\pi^*\cF^H}\big)$ is a $G$-equivariant quasi-coherent sheaf. Since the restriction of the isomorphism $\uptheta^{\cF}|_{\pi^*\cF^H}\colon {\pi^*\cF^H}\simto {\sigma^*\cF^H}$ to $H\times X$ is the identity, there is a unique $G/H$-equivariant structure $\uptheta^{\cF^H}\colon \overline{\pi}^*\cF^H\simto \overline{\sigma}^*\cF^H$ on $\cF^H$ such that $\uptheta^{\cF}|_{\pi^*\cF^H}={(p\times \id _X)}^*\uptheta^{\cF^H}$, where $p\colon G\to G/H$ is the natural projection. This defines the functor
\begin{equation}\label{inv coh}
(-)^H\colon\ \Qcoh_GX\to \Qcoh_{G/H}X.
\end{equation}

If $\uprho\colon\cO_X\to\cA$ is a $G$-equivariant $X$-algebra, the induced morphism
\begin{gather*}
\cO_X=\cO_X^H\xrightarrow{\uprho^H}\cA^H
\end{gather*}
defines a $G/H$-equivariant $X$-algebra. If $\big(\cM,\uptheta^{\cM}\big)$ is a $G$-equivariant $\cA$-module, the morphism $\alpha\colon \cA\times \cM\to \cM$ defining the $\cA$-module structure on $\cM$ is $G$-equivariant since $\uptheta^{\cM}$ is $\pi^*\cA$-linear. The induced morphism
\begin{gather*}
\alpha^H\colon\ \cA^H\times \cM^H\to \cM^H
\end{gather*}
 defines a $\cA^H$-module structure on $\cM^H$, and so the pair $\big(\cM^H,\uptheta^{\cM^H}\big)$ is a $G/H$-equivariant $\cA^H$-module. Thus the functor \eqref{inv coh} extends to the functor
\begin{gather*}
(-)^H\colon\ \Qcoh_G\cA\to \Qcoh_{G/H}\cA^H.
\end{gather*}
The following is a generalization of~\cite[Lemma~2.22]{bfk}.

\begin{prop}
The composition
\[
\begin{tikzcd}
\Qcoh_{G/H}\cA^H\arrow[r,"\id_{p}^*"]&\Qcoh_{G}\cA^H\arrow[rr,"(-)\otimes_{\cA^H}\cA"]&&\Qcoh_{G}\cA
\end{tikzcd}
\]
is left adjoint to the functor $(-)^H$.
\begin{proof}
Let $\cM\in \Qcoh_{G/H}\cA^H$ and $\cN\in \Qcoh_G\cA$. The result follows from the following sequences of isomorphisms
\begin{align*}
\Hom_{\Qcoh_{G}\cA}\big(\id_p^*\cM\otimes_{\cA^H}\cA,\cN\big)
&\cong\Hom_{\Qcoh_{G}\cA^H}\big(\id_p^*\cM,\cN_{\iota}\big)
\\
&\cong \Hom_{\Qcoh_{G/H}\cA^H}\big(\cM,(\cN_{\iota})^H\big)
\\
&\cong \Hom_{\Qcoh_{G/H}\cA^H}\big(\cM,\cN^H\big),
\end{align*}
\looseness=1
where $\cN_{\iota}$ is the restriction of $\cN$ by the natural inclusion $\iota\colon \cA^H\hookto \cA$, the first isomorphism follows from \eqref{rest ext adj}, the second isomorphism follows from an identical argument as in the proof of~\cite[Lemma~2.22]{bfk}, and the last isomorphism follows from a natural isomorphism \mbox{$(\cN_{\iota})^H\cong \cN^H$}.
\end{proof}
\end{prop}

\subsubsection{Restriction functors and induction functors}
 Let $H$ be a closed subgroup of $G$, and
 \begin{gather*}
 \alpha\colon\ H\hookto G
 \end{gather*}
 the natural inclusion morphism. Let $\cA$ be a $G$-equivariant $X$-algebra. We define the {\it restriction functor}
\begin{gather*}
\Res^G_H\defeq (\id_X)_{\alpha}^*\colon\ \Qcoh_G\cA\to \Qcoh_H\cA
\end{gather*}
to be the pull-back by the identity morphism $\id_X\colon X\to X$ that is $\alpha$-equivariant. If $H$ is trivial, we write $\Res\defeq\Res^G_{\{1\}}$, which is nothing but the forgetful functor $\big(\cM,\uptheta^{\cM}\big)\mapsto \cM$.

Next, we will construct the adjoint functor of this restriction functor. We define an $H$-action on $G\times X$ by
\begin{gather*}
h\cdot(g,x)\defeq\big(g\big(h^{-1}\big), h x\big)
\end{gather*}
for any $h\in H$, $g\in G$ and $x\in X$, and we write
\begin{gather*}
G\timesH X\defeq[G\times X/H]
\end{gather*}
for the associated quotient stack. By~\cite[Lemma~2.16(a)]{bfk}, the quotient stack $G\timesH X$ is representable by the scheme $G/H\times X$.
Since the morphism $\sigma\colon G\times X\to X$ defining the $G$-action on $X$ is $H$-invariant, we have the induced morphism
\begin{gather*}
\sigma^H\colon\ G\timesH X\to X.
\end{gather*}
We define the morphism
\begin{gather*}
\varepsilon_X^H\colon\ X\to G\timesH X
\end{gather*}
to be the composition $X\xrightarrow{\varepsilon_X}G\times X\xrightarrow{\,\,q\,\,}G\timesH X$,
where $q$ is the canonical quotient map. Then we have $\sigma^{H}\circ\varepsilon_X^H=\id_X$. The $G$-action on $G\times X$ given by
\begin{gather}\label{3b}
\widetilde{\sigma}\colon\ G\times (G\times X)\to G\times X,\qquad
(g,g',x)\mapsto (gg',x)
\end{gather}
induces the $G$-action on $G\timesH X$. With respect to this $G$-action on $G\timesH X$, the morphism $\varepsilon_X^H$ is~$\alpha$-equivariant, and $\sigma^H$ is $G$-equivariant. Thus $\big(\sigma^H\big)^*\cA$ is a $G$-equivariant $G\timesH X$-algebra, and we have the following functor
\begin{gather*}
\Phi\defeq\big(\varepsilon_X^H\big)^*_{\alpha}\colon\ \Qcoh _G\big(\sigma^H\big)^*\cA\to \Qcoh_H\cA.
\end{gather*}
The following is a generalization of~\cite[Lemma~2.13]{bfk}

\begin{lem}%\label{epsilon equiv}
The functor $\Phi\colon \Qcoh _G\big(\sigma^H\big)^*\cA\to \Qcoh_H\cA$ is an equivalence.
\end{lem}
\begin{proof} This can be proved by a similar argument as in~\cite[Lemma~1.3]{tho}.
 Since $\big(\sigma^H\big)^*\cA$ is isomorphic to the restriction of the big fppf sheaf $[\sigma^*\cA/H]$ to the small Zariski site of~\mbox{$G\timesH X$}, by Proposition~\ref{quotient correspond} we have an equivalence $\Qcoh\big(\sigma^H\big)^*\cA\cong \Qcoh_H\sigma^*\cA$. Thus objects in $\Qcoh _G\big(\sigma^H\big)^*\cA$ are identified with pairs
\begin{equation}\label{object 1}
\Big(\big(\cM,\uptheta^{\cM}\big), \uptheta^{\left(\cM,\uptheta^{\cM}\right)}\Big)
\end{equation}
{\sloppy of $\big(\cM,\uptheta^{\cM}\big)\in \Qcoh_H\sigma^*\cA$ and an isomorphism $\uptheta^{\left(\cM,\uptheta^{\cM}\right)}\colon \widetilde{\pi}^*\big(\cM,\uptheta^{\cM}\big)\simto \widetilde{\sigma}^*\big(\cM,\uptheta^{\cM}\big)$ in $\Qcoh_H(\sigma\circ\widetilde{\sigma})^*\cA$ satisfying the conditions as in \eqref{compati}, where $\widetilde{\pi}\colon G\times G\times X\to G\times X$ and $\widetilde{\sigma}\colon G\times G\times X\to G\times X$ are the projection and the group action in \eqref{3b} respectively. If we write $\widetilde{\uptheta}^{\cM}\colon \widetilde{\pi}^*\cM\simto\widetilde{\sigma}^*\cM$ for the isomorphism in $\Qcoh(\sigma\circ\widetilde{\sigma})^*\cA$ defining the isomorphism $\uptheta^{\left(\cM,\uptheta^{\cM}\right)}$, then the pair $(\cM,\widetilde{\uptheta}^{\cM})$ is an object in $\Qcoh_G\sigma^*\cA$, and $\uptheta^{\cM}$ defines a $H$-equivariant structure on~$\big(\cM,\widetilde{\uptheta}^{\cM}\big)$. Thus the object~\eqref{object 1} uniquely corresponds to an object in $\Qcoh_H[\sigma^*\cA/G]$. Since $G$ trivially acts on $X$ in the $G$-action~\eqref{3b}, the composition of $\varepsilon_X\colon X\to G\times X$ and the natural projection $G\times X\to [G\times X/G]$ gives an isomorphism $ X\simto[G\times X/G]$. Via this isomorphism, the sheaf $[\sigma^*\cA/G]$ on $[G\times X/G]$ corresponds to the sheaf $\cA$ on $X$. Hence the assignment
 \begin{gather*}
 \Bigl(\big(\cM,\uptheta^{\cM}\big), \uptheta^{\left(\cM,\uptheta^{\cM}\right)}\Bigr)\mapsto \big(\varepsilon_X^*\cM,\uptheta^{\cM}|_{H\times \{1\}\times X}\big)
 \end{gather*}}\noindent
 defines an equivalence $\Qcoh_G[\sigma^*\cA/H]\simto \Qcoh_H\cA$, and $\Phi$ is isomorphic to this equivalence via $\Qcoh _G\big(\sigma^H\big)^*\cA\cong \Qcoh_G[\sigma^*\cA/H]$.
 \end{proof}


Since the morphism $\sigma^H\colon G\timesH X\to X$ is $G$-equivariant, we have the direct image functor $\sigma^H_*\colon\Qcoh_G\big(\sigma^{H}\big)^*\cA\to\Qcoh_G\cA$, and
we define the {\it induction functor}
\begin{gather*}
\Ind^G_H\colon\ \Qcoh_H \cA\to \Qcoh_G\cA
\end{gather*}
to be the composition $\big(\sigma^H\big)_*\circ \Phi^{-1}$. If $H$ is trivial, we write $\Ind\defeq\Ind^G_{\{1\}}$.
\begin{lem}\label{split mono} Notation is the same as above.
\begin{itemize}\itemsep=0pt
\item[$1.$] The functor $\Res^G_H$ is exact, and if $G/H$ is affine, the functor $\Ind^G_H$ is also exact.
\item[$2.$] The restriction functor $\Res^G_H$ is left adjoint to the induction functor $\Ind^G_H$.
\item[$3.$] If $H$ is a normal subgroup of $G$ and $G/H$ is reductive affine algebraic group, the adjunction morphism
\begin{gather*}
\mu\colon\ \id\to \Ind^G_H\circ \Res^G_H
\end{gather*}
is a split mono, i.e., there exists a functor morphism $\nu\colon\Ind^G_H\circ\Res^G_H\to \id$ such that the composition $\nu\circ\mu$ is the identity morphism of functors.

\end{itemize}
\begin{proof}
(1) Since the morphism $\sigma^H$ is flat, the pull-back $\big(\sigma^H\big)^*$ is exact. Since $\sigma^{H}\circ\varepsilon_X^H=\id_X$, the functor $\Res^G_H$ is isomorphic to the composition $\Phi\circ\big(\sigma^{H}\big)^*$, and thus $\Res^G_H$ is exact. If $G/H$ is affine, the morphism $\sigma^H$ is an affine morphism since $G\timesH X$ is isomorphic to $G/H\times X$. Then the direct image $\big(\sigma^{H}\big)_*$ is exact, and so is $\Ind^G_H$.

(2) This follows from the adjunction $\big(\sigma^{H}\big)^*\dashv\big(\sigma^{H}\big)_*$ and the isomorphism $\Res^G_H\cong \Phi\circ\big(\sigma^{H}\big)^*$.

(3) It is enough to show that the adjunction morphism
\begin{gather*}
\widetilde{\mu}\colon\ \id \to \big(\sigma^H\big)_*\big(\sigma^H\big)^*
\end{gather*}
is a split mono.
 Note that we have the following cartesian square
\[
\begin{tikzcd}
G\timesH X\arrow[rr,"\sigma^H"]\arrow[d]&&X\arrow[d, "q"]
\\
G/H\arrow[rr, "p"]&&\Spec k,
\end{tikzcd}
\]
where $p$ and $q$ are the morphisms defining the base field $k$.
If $G/H$ is a reductive affine algebraic group, it is linearly reductive. Then the natural morphism
$k\to p_*p^*k=k[G/H]$
is a split mono, and so is the adjunction morphism
\begin{gather*}
\cO_X\to\big(\sigma^H\big)_*\big(\sigma^H\big)^*\cO_{X}
\end{gather*}
by the base change formula for the above cartesian square. For any $\cM\in \Qcoh_G\cA$, by~Lem\-ma~\ref{proj form} the adjunction $\widetilde{\mu}(\cM)\colon \cM\to \big(\sigma^H\big)_*\big(\sigma^H\big)^*\cM$ is isomorphic to the tensor product \begin{gather*}\Bigr(\cO_X\to\big(\sigma^H\big)_*\big(\sigma^H\big)^*\cO_{X}\Bigr)\otimes_{\cO_X}\cM,\end{gather*} and therefore $\widetilde{\mu}(\cM)$ is also a split mono.
\end{proof}
\end{lem}


We will apply the above result to the following:

\begin{lem}\label{faithful general}
Let $\scrC$ and $\scrD$ be additive categories, and
$F\colon \scrC\to \scrD$ and $G\colon \scrD\to \scrC$ additive functors. If $F$ is left adjoint to $G$, and if the adjunction morphism $\mu\colon \id_{\scrC}\to G\circ F$ is a split mono, then the functor $F$ is faithful.
\begin{proof}
For a morphism $f\colon A\to B$ in $\scrC$, assume that $F(f)=0$. It is enough to show that $f=0$. We have a commutative diagram
\[
\begin{tikzcd}
A\arrow[rr, "\mu(A)"]\arrow[rrrr, bend left, "\id"]\arrow[d, "f"']&&GF(A)\arrow[d, "GF(f)"]\arrow[rr, "\nu(A)"]&&A\arrow[d, "f"]\\
B\arrow[rr, "\mu(B)"]&&GF(B)\arrow[rr, "\nu(B)"]&& B.
\end{tikzcd}
\]
Therefore, $f=\nu(B)\circ GF(f)\circ \mu(A)=\nu(B)\circ G(0)\circ \mu(A)=0$.
\end{proof}
\end{lem}




\subsection{Fundamental properties of equivariant modules} In this subsection, we prove that the category of equivariant modules over a certain equivariant algebra has enough injectives and enough locally free modules.

Let $X$ be a quasi-compact and quasi-separated scheme with an action from an affine algebraic group $G$, and $\uprho\colon \cO_X\to\cA$ a $G$-equivariant $X$-algebra.

\begin{prop}\label{grothen prop}
 The category $\Qcoh_G\cA$ is a Grothendieck category. In particular, it has enough injectives.
\begin{proof}
The category $\Qcoh_G\cA$ is an abelian category with small direct sums, and so it suffices to prove that (1) filtered colimits are exact in $\Qcoh_G\cA$ and that (2) $\Qcoh_G\cA$ has a generator.

(1) For a point $x\in X$, denote by $F_x\colon \Qcoh_G\cA\to \Mod\cA_x$ the composition
\begin{gather*}
\Qcoh_G\cA\xrightarrow{\Res}\Qcoh\cA\xrightarrow{(-)_x}\Mod\cA_x,
\end{gather*}
where the latter is taking the stalk at $x$, and $\cA_x$ is the stalk of $\cA$ at $x$. Then a sequence in~$\Qcoh_G\cA$ is exact if and only if for every $x\in X$ the sequence in $\Mod\cA_x$ induced by $F_x$ is exact. Therefore, since $F_x$ commutes with filtered colimits and $\Mod \cA_x$ is a Grothendieck category, filtered colimits in $\Qcoh_G\cA$ are also exact.

(2)
It is well known that $\Qcoh_GX$ is a Grothendieck category (more generally, the category of quasi-coherent sheaves on an algebraic stack is a Grothendieck category~\cite[Proposition~14.2; 0781]{stacks}, \cite[Corollary~5.10]{tlrv}). In particular, $\Qcoh_GX$ has a generator $\cG\in \Qcoh_G X$. We show that
\begin{gather*}\cG_{\cA}\defeq\cG\otimes_{\cO_X}\cA\in \Qcoh_G \cA\end{gather*} is a generator. Let $\cM\in\Qcoh_G\cA$ be an object. Then, there is a set $I$ and a surjective morphism
\begin{gather*}
p\colon\ \cG^{\oplus I}\surjto \cM_{\uprho}
\end{gather*}
in $\Qcoh_G X$. Since the functor $(-)\otimes_{\cO_X}\cA\colon \Qcoh_G X\to \Qcoh_G\cA$ is left adjoint to the functor $(-)_{\uprho}$, it is right exact and commutes with small direct sums. Hence the extension of $p$ by $\uprho\colon \cO_X\to \cA$ defines a surjective morphism
\begin{gather*}
p_{\cA}\colon\ \cG_{\cA}^{\oplus I}\surjto \cM_{\uprho}\otimes_{\cO_X}\cA.
\end{gather*}
Composing this with a natural surjection $\cM_{\uprho}\otimes_{\cO_X}\cA\to \cM$, we have a surjective morphism $\cG_{\cA}^{\oplus I}\surjto \cM$. This shows that $\cG_{\cA}$ is a generator in $\Qcoh_G \cA$.
\end{proof}
\end{prop}





{\samepage\begin{dfn} Notation is the same as above.
\begin{itemize}\itemsep=0pt
\item[1.] A $G$-equivariant $\cA$-module $\big(\cM,\uptheta^{\cM}\big)$ is said to be {\it locally free}, if there is an open covering $\{U_i\hookto X\}_{i\in I}$ of $X$ such that for every $U_i$ the restriction $\cM|_{U_i}$ of the underlying $\cA$-module~$\cM$ is a free $\cA|_{U_i}$-module.

\item[2.] We say that $\big(\cA,\uptheta^{\cA}\big)$ {\it satisfies the resolution property}, if for any $G$-equivariant coherent $\cA$-module $\cN\in \coh_G\cA$, there exists a surjective morphism $\big(\cE,\uptheta^{\cE}\big)\surjto \big(\cN,\uptheta^{\cN}\big)$ in $\coh_G\cA$ from a $G$-equivariant locally free coherent $\cA$-module $\big(\cE,\uptheta^{\cE}\big)$.
\item[3.] A $G$-scheme $X$ {\it satisfies the $G$-equivariant resolution property} if the $G$-equivariant $X$-al\-ge\-bra $(\cO_X,\uptheta^{\can})$ has the resolution property.
\end{itemize}
\end{dfn}

}

\begin{rem}
(1) If $X=\Spec R$ is an affine noetherian scheme, a $G$-equivariant coherent $R$-algebra $A$ satisfies the resolution property if and only if the category $\Mod_GA$ has enough projectives.

(2) Assume that $X$ is noetherian and that $\cA$ is coherent. If $\big(\cA,\uptheta^{\cA}\big)$ satisfies the resolution property, for every $G$-equivariant $\cA$-module $\cM$, there is a surjection $\cE\surjto \cM$ from a $G$-equivariant locally free $\cA$-module $\cE$. This follows since there is a set $\{\cM_i\}_{i\in I}$ of $G$-equivariant coherent submodules $\cM_i$ of $\cM$ such that the natural map $\bigoplus_{i\in I}\cM_i\to \cM$ is surjective.
\end{rem}


\begin{prop}\label{resol prop}
Assume that $\big(\cA,\uptheta^{\cA}\big)$ is coherent and that $X$ has the $G$-equivariant resolution property. Then $\big(\cA,\uptheta^{\cA}\big)$ satisfies the resolution property.
\begin{proof}
Let $\cM\in \coh_G\cA$ be a $G$-equivariant coherent $\cA$-module. By the assumption, there is a surjective morphism $p\colon \cE\surjto \cM_{\uprho}$ from a $G$-equivariant locally free coherent $\cO_X$-module $\cE$. Then the extension
\begin{gather*}
p\otimes_{\cO_X}\cA\colon\ \cE\otimes _{\cO_X}\cA\to \cM_{\uprho}\otimes_{\cO_X}\cA
\end{gather*}
of $p$ by $\uprho\colon \cO_{X}\to \cA$ is also a surjection in $\Qcoh_G \cA$. Composing with the natural surjection $\cM_{\uprho}\otimes_{\cO_{X}} \cA\surjto \cM$ gives a surjective morphism $\cE\otimes_{\cO_X} \cA\surjto \cM$ from a $G$-equivariant locally free $\cA$-module.
\end{proof}
\end{prop}

\begin{dfn}
A $G$-equivariant (resp.\ coherent) $\cA$-module $\big(\cM,\uptheta^{\cM}\big)$ has a {\it finite locally free resolution} in $\Qcoh_G\cA$ (resp.\ in $\coh_G\cA$) if there exists an exact sequence
\begin{gather*}
0\to \cE^{-n}\to\cdots\to\cE^0\to \cM\to0
\end{gather*}
in $\Qcoh_G\cA$ (resp.\ in $\coh_G\cA$) such that each $\cE^i$ is locally free.
\end{dfn}


\begin{lem}\label{enough proj}
Assume that $X$ is a normal scheme of finite type over $k$ and that $G$ is affine.
If~$X$ has an ample family of line bundles, then $\big(\cA,\uptheta^{\cA}\big)$ satisfies the resolution property.
\begin{proof}
This follows from~\cite[Theorem~2.29]{bfk} and Proposition~\ref{resol prop}.
\end{proof}
\end{lem}

\begin{lem}\label{finite proj resol}
Let $R$ be a normal ring of finite type over $k$, and $G$ an affine reductive algebraic group acting on $\Spec R$. Let $\big(A,\uptheta^{A}\big)$ be a $G$-equivariant coherent $R$-algebra, and $\big(M,\uptheta^M\big)\in \Mod_GA$ a $G$-equivariant $A$-module. If the underlying $A$-module $M\in \Mod A$ has a finite projective resolution, so does $\big(M,\uptheta^M\big)$. If $M$ is finitely generated and has a finite projective resolution in $\fmod A$, then $\big(M,\uptheta^M\big)$ has a finite projective resolution in $\fmod_GA$.
\begin{proof}
By Lemma~\ref{enough proj}, the categories $\Mod_GA$ and $\fmod_GA$ have enough projectives. Since~$A$ is coherent $R$-algebra, $\fmod_GA$ is an abelian subcategory of $\Mod_GA$. It is standard that an~object~$x$ in an abelian category $\scrA$ has a finite projective resolution if and only if there exists $n>0$ such that $\Ext_{\scrA}^i(x,y)=0$ for any $i>n$ and any $y\in \scrA$. Thus the statements follow from Proposition~\ref{G hom 2}.
\end{proof}
\end{lem}

\subsection{Equivariant tilting modules and derived equivalences}\label{derived section}

In this subsection, we prove that an equivariant tilting module induces a derived equivalence of~equivariant algebras, which is considered as an equivariant Morita theory.

Let $X$ be a separated scheme of finite type over $k$, and $G$ an affine algebraic group acting on~$X$.
Let $H\subseteq G$ be a closed normal subgroup of $G$, and $\cA$ a $G$-equivariant coherent $X$-algebra.

\begin{dfn}
Let $\cT\in \coh_G\cA$ be a $G$-equivariant coherent $\cA$-module.
\begin{itemize}\itemsep=0pt
\item[1.] $\cT$ is called a {\it $(G,H)$-tilting module} (resp.\ {\it partial $(G,H)$-tilting module}), if the restriction $\Res^G_H(\cT)$ is a tilting object (resp.\ partial tilting object) in $\Qcoh_H\cA$.
\item[2.] $\cT$ is called a {\it $G$-tilting module} (resp.\ {\it partial $G$-tilting module}) if it is $(G,\{1\})$-tilting (resp.\ partial $(G,\{1\})$-tilting).
\end{itemize}
\end{dfn}
Let $R$ be a normal ring of finite type over $k$, and suppose that there is a $G$-action $\sigma\colon G\times \Spec R\to \Spec R$ on $\Spec R$ such that the $H$-action on $\Spec R$, which is given by the restriction of $\sigma$, is trivial. Let $f\colon X\to \Spec R$ be a $G$-equivariant morphism such that $f_*\cO_X=R$ and the associated morphism $\overline{f}\colon [X/H]\to \Spec R$ is proper. Let $\big(\cT,\uptheta^{\cT}\big)\in \coh_G\cA$ be
a $G$-equivariant coherent $\cA$-module. Then the endomorphism ring
\begin{gather*}
\Lambda\defeq\End_{\cA}(\cT)=f_*\cEnd_{\cA}(\cT)
\end{gather*}
of the underlying coherent $\cA$-module $\cT$ is a $G$-equivariant $R$-algebra.
We define the functor
\begin{gather*}
\F\colon\ \Qcoh_G\cA\to \Mod_{G/H}\Lambda^H
\end{gather*}
to be the composition
\[
\begin{tikzcd}
\Qcoh_G\cA\arrow[rr,"\cHom_{\cA}(\cT\mbox{\hspace{-0.5mm}\tiny ,}-)"]&&\Qcoh_G\cEnd_{\cA}(\cT)\arrow[r,"f_{\star}"]&\Mod_G\Lambda\arrow[r,"(-)^H"]&\Mod_{G/H}\Lambda^H.
\end{tikzcd}
\]
Then the composition
\begin{gather*}
\G\defeq\left(-\otimes \cT\right)\circ f^{\star}\circ \left(\id_p^*(-)\otimes_{\cA^H}\cA\right) \colon\ \Mod_{G/H}\Lambda^H\to \Qcoh_G\cA
\end{gather*}
is left adjoint to $\F$, and it preserves equivariant coherent modules. Note that since the morphism $\overline{f}\colon [X/H]\to \Spec R$ is proper, the functor $\F$ also preserves equivariant coherent modules. The following is one of our motivations of considering equivariant algebras and equivariant modules.

\begin{thm}\label{derived equiv}
Assume that $G/H$ is reductive and that every $\Lambda^H$-module $M\in \Mod\Lambda^H$ is of finite projective dimension.
\begin{itemize}\itemsep=0pt
\item[$1.$]
If $\big(\cT,\uptheta^{\cT}\big)$ is partial $(G,H)$-tilting, then the functor
\begin{gather*}
\bL \G\colon\ \Db\big(\Mod_{G/H}\Lambda^H\big)\hookto \Db\big(\Qcoh_G{\cA}\big)
\end{gather*}
is fully faithful, and it restricts to the fully faithful functor
\begin{gather*}
\bL \G\colon\ \Db\big(\fmod_{G/H}\Lambda^H\big)\hookto \Db\big(\coh_G{\cA}\big).
\end{gather*}
\item[$2.$] If $\big(\cT,\uptheta^{\cT}\big)$ is $(G,H)$-tilting, then the functor
\begin{gather*}
\bL \G\colon\ \Db\big(\Mod_{G/H}\Lambda^H\big)\simto \Db\big(\Qcoh_G{\cA}\big)
\end{gather*}
is an equivalence, and it restricts to the equivalence
\begin{gather*}
\bL \G\colon\ \Db\big(\fmod_{G/H}\Lambda^H\big)\simto \Db\big(\coh_G{\cA}\big).
\end{gather*}
\end{itemize}
\end{thm}

Before we prove this theorem, we prepare the following lemmas.
\begin{lem}\label{faithful lemma}
Notation is the same as above, and assume that $G/H$ is reductive.
Then the restrictions
\begin{gather*}
\Res^G_H\colon\ \Db\big(\Qcoh_G\cA\big)\to \Db\big(\Qcoh_H\cA\big),
\\[.5ex]
\Res\colon\ \Db\big(\Mod_{G/H}\Lambda^H\big)\to \Db\big(\Mod\Lambda^H\big)
\end{gather*}
are faithful functors.
\begin{proof}
We only prove $\Res^G_H$ is faithful. Since $G/H$ is affine, $\Ind^G_H\colon \Qcoh_H\cA\to \Qcoh_G\cA$ is an~exact functor, and so it extends to the functor
\begin{gather*}
\Ind^G_H\colon\ \Db\big(\Qcoh_H\cA\big)\to \Db\big(\Qcoh_G\cA\big)
\end{gather*}
which is right adjoint to $\Res^G_H$.
 Since $G/H$ is reductive, the adjunction $\mu\colon \id_{\Qcoh_G\cA} \to\Ind^G_H\Res^G_H$ is a split mono by Lemma~\ref{split mono}, and this splitting naturally extends to the splitting of the adjunction $\mu\colon \id_{\Db\left(\Qcoh_G\cA\right)} \to\Ind^G_H\Res^G_H$ since $\Ind^G_H$ and $\Res^G_H$ are exact functors. Hence the functor $\Res^G_H$ is faithful by Lemma~\ref{faithful general}.
\end{proof}
\end{lem}

\begin{lem}\label{reduction lemma}
For each $i=1,2$, let $F_i\colon \scrC_i\to \scrD_i$ and $G_i\colon \scrD_i\to \scrC_i$ be exact functors between triangulated categories such that $G_i\dashv F_i$. Denote by $\sigma_i\colon G_i\circ F_i\to \id$ and $\tau_i\colon \id\to F_i\circ G_i$ the adjunction morphisms, and let $P\colon \scrC_1\to \scrC_2$ and $Q\colon \scrD_1\to\scrD_2$ be faithful exact functors. Assume that we have functor isomorphisms $\varphi\colon Q\circ F_1\simto F_2\circ P$ and $\psi\colon P\circ G_1\simto G_2\circ Q$ and that the following diagrams are commutative
\[
\begin{tikzcd}
PG_1F_1\arrow[rr, "\sim"]\arrow[rd, "P\sigma_1"']&&G_2F_2P\arrow[ld, "\sigma_2P"]&&
QF_1G_1\arrow[rr, "\sim"]&&F_2G_2Q\\
&P&&&&Q\arrow[lu, "Q\tau_1"]\arrow[ru, "\tau_2Q"']&
\end{tikzcd}
\]
where the top horizontal arrows are the isomorphisms induced by $\varphi$ and $\psi$.
\begin{itemize}\itemsep=0pt
\item[$1.$] If $F_2$ $($resp.\ $G_2)$ is fully faithful, so is $F_1$ $($resp.\ $G_1)$.
\item[$2.$] If $F_2$ $($resp.\ $G_2)$ is an equivalence, so is $F_1$ $($resp.\ $G_1)$.
\end{itemize}
\begin{proof}
(2) follows from (1). We only prove (1) for $F_i$, since the statement for $G_i$ follows by a~simi\-lar argument. For an object $C\in \scrC_1$ consider the following triangle
\begin{gather*}
 G_1F_1(C)\xrightarrow{\sigma_1(C)} C\to{\rm Cone}(\sigma_1(C))\to G_1F_1(C)[1].
\end{gather*}
It suffices to prove that ${\rm Cone}(\tau_1(C))$ is the zero object. By the assumption, $P({\rm Cone}(\sigma_1(C)))$ is isomorphic to the cone of the morphism $\sigma_2(P(C))$, which is the zero object since $F_2$ is fully faithful. Hence ${\rm Cone}(\sigma_1(C))$ is also the zero object since $P$ is faithful.
\end{proof}
\end{lem}




\noindent
{\bf Proof of Theorem~\ref{derived equiv}.}
We prove (1) and (2) simultaneously. We define a functor
\begin{gather*}
\F_H\colon\ \Qcoh_H\cA\to \Mod\Lambda^H
\end{gather*}
to be the composition
\[
\begin{tikzcd}
\Qcoh_H\cA\arrow[rr,"\cHom_{\cA}(\cT\mbox{\hspace{-0.5mm}\tiny ,}-)"]&&\Qcoh_H\cEnd_{\cA}(\cT)\arrow[r,"f_{\star}"]&\Mod_H\Lambda\arrow[r,"(-)^H"]&\Mod\Lambda^H.
\end{tikzcd}
\]
If we set $\cT_H\defeq\Res^G_H\big(\cT,\uptheta^{\cT}\big)\in \Qcoh_H\cA$, since $\Lambda^H\cong \End_{\cA}(\cT)^H=\End_{\Qcoh_H\cA}(\cT_H)$, the functor $\F_H$ is isomorphic to the functor
\begin{gather*}
\Hom_{\Qcoh_H\cA}(\cT_H,-)\colon\ \Qcoh_H\cA\to \Mod\Lambda^H.
\end{gather*}
Note that $\F_H$ has a left adjoint functor $\G_H\colon \Mod\Lambda^H\to \Qcoh_H\cA$ which is given by
\begin{gather*}
\G_H(-)\defeq f^{\star}\big(\id_p^*(-)\otimes_{\Lambda^H}\Lambda\big)\otimes_{\cEnd(\cT)}\cT\cong f^{-1}\big(\id_p^*(-)\otimes_{\Lambda^H}\Lambda\big)\otimes_{f^{-1}\Lambda}\cT,
\end{gather*}
where $p\colon H\to H/H=\{1\}$ is the natural projection. By construction, we have natural isomorphisms $\F_H(\cT_H)\cong \Lambda^H$ and $\G_H\big(\Lambda^H\big)\cong \cT_H$, and this implies that the adjunction morphisms
\begin{gather*}
\Lambda^H\to \F_H\G_H\big(\Lambda^H\big) \qquad\text{and}\qquad \G_H\F_H(\cT_H)\to \cT_H
\end{gather*}
are isomorphisms.
Since $\cT_H$ is a partial tilting object in $\Qcoh_H\cA$, by Theorem~\ref{general tilting}, the functor
\begin{gather*}
\bL\G_H\colon\ \Db\big(\Mod\Lambda^H\big)\to \Db\big(\Qcoh_H\cA\big)
\end{gather*}
is fully faithful, and if $\cT_H$ is tilting, it is an equivalence.
Now we have the following commutative diagram:
\[
\begin{tikzcd}
\Db\big(\Mod_{G/H}\Lambda^H\big)\arrow[rr,"\bL\G"]\arrow[d, "\Res"']&&\Db\big(\Qcoh_G\cA\big)\arrow[d, "\Res^G_H"]\arrow[rr,"\bR\F"]&&\Db\big(\Mod_{G/H}\Lambda^H\big)\arrow[d, "\Res"']\\
\Db\big(\Mod\Lambda^H\big)\arrow[rr, "\bL \G_H"]&&\Db\big(\Qcoh_H\cA\big)\arrow[rr, "\bR \F_H"]&&\Db\big(\Mod\Lambda^H\big).
\end{tikzcd}
\]
Hence (1) and (2) follows from Lemmas~\ref{faithful lemma} and~\ref{reduction lemma}.\qed
